\chapter{Penis Irritation (Balanitis)}

\section*{Definition}
Inflammation or irritation of the glans penis and/or foreskin (balanoposthitis), characterized by redness, swelling, itching, discharge, and discomfort.

\section*{What Happens in Your Body}
Irritation results from infectious (candidal, bacterial, viral, parasitic), inflammatory (lichen sclerosus, contact dermatitis, psoriasis), or traumatic (friction, chemical irritation) causes. Poor hygiene, diabetes, and immunosuppression predispose to infection. Phimosis (tight foreskin) traps secretions increasing irritation risk.

\section*{Causes}
\begin{itemize}
  \item Candidal infection (yeast — especially in diabetes)
  \item Poor hygiene (smegma accumulation)
  \item Contact dermatitis (soaps, detergents, lubricants, latex condoms)
  \item Bacterial infection (streptococcal, staphylococcal, anaerobic)
  \item Sexually transmitted infections (HSV, gonorrhea, chlamydia, syphilis)
  \item Lichen sclerosus (chronic inflammatory skin condition)
  \item Psoriasis or reactive arthritis
  \item Trauma or friction (vigorous intercourse, masturbation)
  \item Phimosis or Tight foreskin
  \item Allergic reaction (medications, topical creams)
\end{itemize}

\section*{When to See a Doctor}
See your doctor if:
\begin{itemize}
  \item Symptoms are severe or getting worse over time
  \item Penis irritation with discharge, fever, difficulty urinating, or inability to retract foreskin (paraphimosis)
  \item You have other concerning symptoms such as fever, weight loss, or bleeding
\end{itemize}

\section*{Related Symptoms}
\begin{itemize}
  \item Penile discharge
  \item Dysuria
  \item Swelling
  \item Itching
  \item Painful urination
\end{itemize}
\textbf{Important:} If this symptom persists, your doctor will check for serious underlying causes.

\section*{Self-Care / Home Management}
\begin{itemize}
  \item Rest and avoid activities that make it worse.
  \item Maintain adequate hydration and a balanced diet.
  \item Over-the-counter remedies may provide symptomatic relief (consult a pharmacist or doctor).
  \item Monitor symptoms and seek professional advice if they persist or worsen.
  \item Avoid self-medicating with prescription medications without medical guidance.
\end{itemize}

\section*{How Your Doctor Will Evaluate You}
\begin{itemize}
  \item A thorough history and physical examination are the first steps in evaluation.
  \item Laboratory tests (blood work, urinalysis) may be ordered based on clinical suspicion.
  \item Imaging studies (X-ray, ultrasound, CT, MRI) may be indicated when structural pathology is suspected.
  \item Specialist referral is arranged when the diagnosis remains uncertain or requires specific expertise.
\end{itemize}

\section*{Treatment Options}
\begin{itemize}
  \item Treatment targets the underlying cause identified through diagnostic evaluation.
  \item Supportive care and symptom management are provided while awaiting definitive therapy.
  \item Medications, lifestyle modifications, and physical therapies may be prescribed as appropriate.
  \item Follow-up ensures treatment effectiveness and allows timely adjustments to the care plan.
\end{itemize}

\clearpage